\section{The Plan on a Real Graph}
\label{sec:ch10-the-plan-on-a-real-graph}
\index{query plan!on Mosaic|(}%

Chapter~\ref{ch:how-a-federated-query-actually-runs} took a query plan apart on
a sample built for the purpose: two subgraphs, three products, two fields. That
is where \code{Sequence} and \code{Parallel}, \code{Single} and
\code{BatchEntity}, \code{dependsOnFetchIds} and the \code{dependencies} array
were established, and none of it is re-derived here. The headers are the same
two as well~\autocite{cosmo2026queryplan}: \code{X-WG-Include-Query-Plan} to get
the plan, and \code{X-WG-Skip-Loader} to get it without executing anything.

What follows is what the sample could not show, because the sample had no
connections, no pagination and one field per subgraph.

Here is the plan for the query in
section~\ref{sec:ch10-the-query-that-stopped-working}. The second fetch carries
one more key, \code{dependencies}, taken out here and printed on its own
further down:

\begin{minted}[fontsize=\footnotesize,breaklines]{json}
{
  "version": "1",
  "kind": "Sequence",
  "children": [
    { "kind": "Single",
      "fetch": {
        "kind": "Single", "subgraphName": "catalog", "subgraphId": "0", "fetchId": 0,
        "query": "query($a: Int){ browseProducts(first: $a){ nodes { title __typename id } } }" } },
    { "kind": "Single",
      "fetch": {
        "kind": "BatchEntity", "path": "browseProducts.nodes",
        "subgraphName": "mosaic", "subgraphId": "1", "fetchId": 1,
        "dependsOnFetchIds": [0],
        "representations": [
          { "kind": "@key", "typeName": "Product",
            "fragment": "fragment Key on Product { __typename id }" } ],
        "query": "query($representations: [_Any!]!, $b: Int){ _entities(representations: $representations){ ... on Product { __typename price { amount currency } availableQuantity averageRating reviews(first: $b){ totalCount nodes { rating author { displayName } } } } } }" } }
  ],
  "normalizedQuery": "query($a: Int, $b: Int){browseProducts(first: $a){nodes {title price {amount currency} availableQuantity averageRating reviews(first: $b){totalCount nodes {rating author {displayName}}}}}}"
}
\end{minted}

\index{listings!whitespace collapsed}%
The router writes those embedded query strings across several lines with
four-space indentation. I have collapsed them, as
chapter~\ref{ch:how-a-federated-query-actually-runs} did, and changed nothing
else.

Two fetches for a document with fourteen selections in it, spanning two
services. The shape is chapter~\ref{ch:how-a-federated-query-actually-runs}'s
exactly. Three details in it are not.

\subsection*{The path points inside the connection}

\index{query plan!path@\code{path}}%
\code{"path": "browseProducts.nodes"}, not \code{"browseProducts"}.

Chapter~\ref{ch:schema-design-that-survives-change} turned \code{reviews} into a
connection and chapter~\ref{ch:data-without-the-n-plus-1} added
\code{browseProducts} as a paged field, so the products are not the root field's
value; they are one level below it, under \code{nodes}. The entity fetch has to
say where in the first fetch's response its representations came from and where
its answers go back, and that is what \code{path} is.

It is a small thing and it is the first sign that Relay-style pagination and
federation have to know about each other. They do, here, and
chapter~\ref{ch:hard-modeling-problems} is where they stop getting along: a
connection whose pages come from one service and whose page contents come from
another is a harder problem than this one.

\subsection*{Four fields, one key, one round trip}

\index{query plan!dependencies array@\code{dependencies} array}%
The \code{dependencies} array I took out has four entries, one for each field
the client asked Mosaic for. Flattened, and with the
\code{isUserRequested: true} that every one of them carries left off:

\begin{minted}[fontsize=\footnotesize]{text}
price              dependsOn  catalog.Product.id   isKey: true
availableQuantity  dependsOn  catalog.Product.id   isKey: true
averageRating      dependsOn  catalog.Product.id   isKey: true
reviews            dependsOn  catalog.Product.id   isKey: true
\end{minted}

Four entries, one distinct dependency. Every field on the far side of the
boundary needs the same key from the same subgraph, so all four travel in one
\code{BatchEntity} fetch, and the plan has two children rather than five.

\index{query plan!what its length does not tell you}%
The reading lesson is in there, and it took me a second look to get right. The
length of that array is not a count of round trips and not a count of anything
expensive. It is one entry per user-requested field, and a chatty-looking array
can be a single fetch. The number that costs you something is the number of
\code{fetch} nodes, and the thing that determines it is how many distinct
dependencies there are, not how many fields have one.

Every entry here says the same thing, because every field crosses on one key.
\code{@requires} is what should make one of them say something else, and
chapter~\ref{ch:entity-resolution-done-right} is where that gets measured
rather than predicted. Mosaic has no \code{@requires} yet, so this chapter has
not seen a non-uniform \code{dependencies} array and does not claim to know
what one costs.

\subsection*{The router rewrote the question}

\index{query plan!normalisation of literals}%
Look at what fetch 0 sends:

\begin{minted}{graphql}
query($a: Int){ browseProducts(first: $a){ nodes { title __typename id } } }
\end{minted}

The client sent \code{first: 3}. A literal. What went to Catalog is a variable,
and the plan's \code{normalizedQuery} shows the same thing happening to the
client's own document, which now opens \code{query(\$a: Int, \$b: Int)} with
both numbers lifted out.

The \code{\_\_typename} and \code{id} in that fetch are
chapter~\ref{ch:how-a-federated-query-actually-runs}'s finding, unchanged: the
router adds them because fetch 1 cannot be built without a key.

The variables are new, and the consequence is worth measuring rather than
inferring. If literals become variables, two clients asking for different page
sizes are sending the router the same document. I checked it three ways.

\medskip\noindent\textbf{Different literals, same plan.} A two-field query at
\code{first: 3}, \code{first: 5} and \code{first: 7} produces three plans that
compare equal as JSON strings, and the storefront query at \code{(3, 2)} and at
\code{(7, 4)} does the same.

\medskip\noindent\textbf{Client-side variables give a different plan.} A client
that declares its variables properly, which is what every generated client
does, sends this:

\begin{minted}{graphql}
query Storefront($first: Int!, $reviews: Int!) {
  browseProducts(first: $first) {
    nodes {
      title
      price { amount currency }
      availableQuantity
      averageRating
      reviews(first: $reviews) {
        totalCount
        nodes { rating author { displayName } }
      }
    }
  }
}
\end{minted}

That normalises to \code{query Storefront(\$a: Int!, \$b: Int!)}, keeping the
operation name and the non-null types. Compared with the inlined form's
\code{query(\$a: Int, \$b: Int)} using \code{===}: false. Two documents that
mean the same thing to a person are two different documents here.

\medskip
\index{query plan!and the plan cache}%
Now the part where I stop. What I have shown is that the plan \emph{documents}
are identical in one case and different in the other. Whether the router keeps a
plan cache, and whether it keys that cache on any of this, I have not looked at,
and chapter~\ref{ch:inside-the-router-and-the-composer} is where the planner
gets read properly. What is fair to take from here is narrower and still useful:
inlining a literal argument costs you nothing at this layer, because the planner
takes it out again, and naming your operation changes the document even though
it does not change the answer.

\subsection*{What the plan still does not tell you}

Everything above is the router's account of what it intends to do, which is
chapter~\ref{ch:how-a-federated-query-actually-runs}'s caution and still holds.
With \code{X-WG-Skip-Loader} set it is an account of what it did not do at all.

For the sample, that chapter closed the gap by logging what the subgraphs
received. Mosaic cannot do that: neither service has request-body logging on,
and section~\ref{sec:ch10-the-query-that-stopped-working} showed that only one
of the two reports anything about a request at all. So for Mosaic the plan is
checked against the answer and against Mosaic's own resolver and SQL counts,
and not against the bytes. That is a weaker check than
chapter~\ref{ch:how-a-federated-query-actually-runs} had, and it is weaker in a
direction the next section makes concrete.

\index{query plan!on Mosaic|)}%
